\section{Bài báo giải được gì, và để lại gì}
\label{sec:ch04-con-lai-gi}

\index{ô nhớ!kết quả trên adding problem}%
Chương~\ref{ch:hien-tuong} huấn luyện một mạng hồi quy thường trên
\tn{bài toán cộng}{adding problem} và xem nó hỏng. Ở \(T = 50\) và \(T = 100\)
thì loss dừng quanh 0.65. Con số ấy không phải một lời giải dở dang: phương sai
của target là 2/3, vì nó là tổng của hai giá trị rút đều từ \([-1, 1]\), nên
một mô hình luôn trả lời 0 đạt đúng 0.6667. Mạng của
chương~\ref{ch:hien-tuong} đạt phương sai. Nó dừng nghe input.

Cần nói ngay một chuyện mà chương~\ref{ch:hien-tuong} đã tự nói và tôi không
được phép quên: trên bài toán này gradient của nó \emph{không} triệt tiêu. Tỉ
số chuẩn gradient giữa bước đầu và bước cuối vẫn quanh 1.0 ngay ở \(T = 100\),
và chương ấy kết luận thẳng rằng bài toán cộng không phải bài kiểm tra tốt cho
\tn{triệt tiêu gradient}{vanishing gradient}. Mạng hỏng, nhưng không hỏng vì
tín hiệu lỗi không đi xa được.

Vậy nên số đo dưới đây không phải bằng chứng cho CEC. Nó là bằng chứng cho
mục~\ref{sec:ch04-hai-xung-dot}. Bài toán cộng đòi đúng cái mà một trọng số
hằng theo thời gian không biểu diễn được: ghi ở hai bước được đánh dấu, làm ngơ
ở chín mươi tám bước còn lại. Một mạng hồi quy thường không có chỗ nào đặt chữ
\enquote{khi nào}, và \(1/T\) của mục~\ref{sec:ch04-hai-xung-dot} là cái giá
của việc thiếu chỗ ấy.

Hai mô hình dưới đây chạy trên cùng bộ sinh, cùng optimizer, cùng số batch và
cùng thứ tự batch với nhau. Chúng không chạy cùng thiết lập với
chương~\ref{ch:hien-tuong}: chương ấy huấn luyện bằng numpy và SGD, còn đây là
PyTorch và Adam. Thứ duy nhất khác nhau giữa hai dòng bảng là cái nằm trong
phép hồi quy.

\begin{measured}{ô nhớ trên bài toán mà chương~\ref{ch:hien-tuong} không giải được}
\begin{minted}{text}
model        d_h  params  loss @400    loss @800    final (mean 100)
LSTM 1997    16   929     0.011706     0.003343     0.001397
plain RNN    32   1153    0.683808     0.672034     0.675564
\end{minted}

\(T = 100\), 1200 batch, batch 64, learning rate 0.01, Adam cho cả hai.

RNN thường được cho tầng ẩn \emph{lớn hơn}, 32 so với 16, nên nó mang 1153
trọng số so với 929 của \tn{ô nhớ}{memory cell}. Kết quả này không đọc được
thành \enquote{LSTM được phát nhiều tham số hơn}.

Ô nhớ về 0.001397, tức 0.002095 phương sai target. RNN thường dừng ở 0.675564,
tức 1.013346 lần phương sai, sau 1200 batch với một optimizer mạnh hơn thứ
chương~\ref{ch:hien-tuong} dùng.

Đo bằng \code{experiments/ch04\_adding.py} tại tag \code{ch04}.
\end{measured}

Cả hai lần huấn luyện gộp lại chạy dưới một phút trên một laptop không có card
đồ họa, và ngân sách mà repo companion tự áp cho mỗi experiment cũng là một
phút. Bài đưa số ở quy mô khác hẳn: quãng vượt 1000 bước, và trên bài toán 2c
của họ thì cả real-time recurrent learning, viết tắt là RTRL, lẫn
\tn{lan truyền ngược qua thời gian}{backpropagation through time}, viết tắt là
BPTT, đều không giải nổi ngay từ quãng 10
bước~\autocite{hochreiter1997lstm}.

\subsection*{Tham số: ba, bốn, hay chín}

\index{ô nhớ!đếm tham số}%
Con số hay được dẫn là bốn: một tầng LSTM có gấp bốn lần trọng số của một tầng
hồi quy thường cùng bề rộng. Bài 1997 thì dẫn một con số khác, \(3^2\). Cả hai
đều đúng và chúng không phải hai câu trả lời cho một câu hỏi.

\begin{minted}{text}
model       blocks  parameters  vs plain RNN
plain RNN   1       131328      1.00
LSTM 1997   3       393984      3.00
LSTM 2000   4       525312      4.00
GRU 2014    3       393984      3.00
\end{minted}

Ba, cho bản 1997 ở dạng tầng: cell input, \tn{cổng vào}{input gate},
\tn{cổng ra}{output gate}. Không có \tn{cổng quên}{forget gate}, nên không có
khối thứ tư. Bốn là bản 2000. Con số bốn nói về một kiến trúc mà chương này còn
chưa tới, và gán nó cho bài 1997 là gán sai.

Còn chín là topology của chính bài, nơi tầng ẩn nối đầy đủ và mỗi cổng nhận kết
nối từ mọi ô nhớ lẫn mọi cổng khác. Thay một hidden unit bằng ba unit thì nhân
ba cả số nguồn lẫn số đích, nên khối hồi quy nhân \(3^2\). Tôi đếm thử ở 64
unit và được đúng 9.0000 trên riêng khối hồi quy; cả tầng thì 8.2603, vì trọng
số input và bias chỉ nhân 3 chứ không nhân 9.

Một chi tiết nhỏ nếu bạn tự đếm bằng PyTorch: \code{torch.nn.LSTM} cho tỷ số
4.0078 chứ không đúng 4, vì nó mang hai vector bias cho mỗi khối thay vì một.

\subsection*{Cái bài tự nêu, và cái chỉ nhìn từ 2026 mới thấy}

\index{ô nhớ!giới hạn}%
Mục 6 của bài liệt kê giới hạn, và danh sách ấy đáng đọc vì nó không phải danh
sách bạn sẽ đoán. Bản cắt cụt không giải được bài không phân rã được, kiểu XOR
trễ. Mỗi \tn{khối ô nhớ}{memory cell block} cần thêm hai unit. Vì dòng lỗi qua
CEC không đổi, LSTM gặp đúng loại khó khăn mà một mạng feedforward nhìn cả
chuỗi một lúc gặp phải, ví dụ bài parity 500 bước. Và mọi phương pháp dựa
gradient đều không đếm chính xác được số bước rời rạc, nên phân biệt
\enquote{99 bước trước} với \enquote{100 bước trước} thì cần thêm một cơ chế
đếm.

Chỗ đáng nói là cái \emph{không} có trong danh sách ấy. Bài không coi tính tuần
tự là một giới hạn. Năm 1997 không ai coi song song hóa là ràng buộc, vì thứ
đang thiếu là khả năng học chứ không phải thông lượng. Người ta hay tóm tắt
khoảng trống mà chương này để lại thành \enquote{vẫn tuần tự, tham số gấp
bốn}, và tôi phải nói rõ: cả hai vế đều là đánh giá nhìn ngược, không phải lời
của bài. Vế thứ hai còn sai một đơn vị so với bản 1997, như bảng trên vừa chỉ
ra.

Nhưng nhìn ngược không có nghĩa là không có thật.
Phương trình~\eqref{eq:ch04-cap-nhat-c} tính \(\vect{c}_t\) từ
\(\vect{c}_{t-1}\), nên \(T\) bước là \(T\) lượt tính nối đuôi nhau, và không
có phần cứng nào rút ngắn được một chuỗi phụ thuộc. Ô nhớ sửa được chuyện
gradient không đi xa; nó không đụng gì tới chuyện phải đi \emph{từng bước một}.
Suốt hai mươi năm sau đó, đó không phải điều đáng phàn nàn. Nó thành điều đáng
phàn nàn đúng lúc kích thước tập dữ liệu và số nhân trên một card vượt qua một
ngưỡng nào đó, và chương~\ref{ch:transformer} là chỗ ai đó quyết định trả giá
để bỏ hẳn phép hồi quy đi.

Còn hai chương nữa trước khi tới đó. Ô nhớ giải bài toán nhớ trong \emph{một}
chuỗi. Câu hỏi tiếp theo không phải nhớ lâu hơn mà là một câu khác hẳn: đọc một
chuỗi rồi sinh ra một chuỗi khác, độ dài khác, thứ tự khác.
Chương~\ref{ch:encoder-decoder} lấy ô nhớ này làm linh kiện và dựng thứ đó, rồi
phát hiện ra \tn{chỗ thắt}{bottleneck} nằm ở một nơi mà chương này không hề
nhắc tới.
